% Generated draft: requires mathematical and source review.
\documentclass[11pt]{article}
\usepackage[margin=25mm]{geometry}
\usepackage{fontspec}
\setmainfont{Latin Modern Roman}
\usepackage{amsmath,amssymb,mathtools,unicode-math}
\setmathfont{Latin Modern Math}
\usepackage{enumitem,needspace,xurl,hyperref}
\setlength{\parindent}{0pt}
\setlength{\parskip}{5pt}
\setlength{\emergencystretch}{3em}
\newcommand{\N}{\mathbb{N}}
\newcommand{\Z}{\mathbb{Z}}
\newcommand{\Q}{\mathbb{Q}}
\newcommand{\R}{\mathbb{R}}
\newcommand{\C}{\mathbb{C}}
\newcommand{\F}{\mathbb{F}}
\newcommand{\A}{\mathbb{A}}
\newcommand{\PP}{\mathbb{P}}
\newcommand{\E}{\mathbb{E}}
\newcommand{\eps}{\varepsilon}
\newcommand{\dd}{\,\mathrm{d}}
\begin{document}
\section*{UnsolvedMath ID: 20001650}
\textbf{Source problem code:} catalogue entry
\subsection{Definitions and mathematical statement}
For a closed connected orientable surface $S_g$, write $\operatorname{Mod}(S_g)$ for the orientation-preserving mapping class group, and let $F_n$ be free of rank $n$. The source's primary question asks to describe the homomorphisms
\[
 \operatorname{Hom}(\operatorname{Mod}(S_g),\operatorname{Out}(F_n)),
\]
and determine when injective ones exist. Its ``closed'' surface assumption does not explicitly fix orientability, connectedness, genus, or extended versus orientation-preserving mapping classes; the displayed convention is the setting of the supplied partial results. The Hom collection is not presumed to be an abelian group.

The source also includes reverse-direction context: it asserts that poison subgroups obstruct embeddings of $\operatorname{Out}(F_n)$ into mapping class groups, citing Brendle and Hamidi-Tehrani; it says an induction argument yields many infinite-image homomorphisms from $\operatorname{Out}(F_3)$ to mapping class groups, and asks about infinite-image maps when $n\geq4$. The poison-subgroup assertion is not given a rank range in the extract and must not be read as an all-rank statement: low-rank exceptions already matter. The aid treats this paragraph as context of Problem 4.2 rather than its primary forward question. Both directions are retained here without assigning a new ID or inventing the omitted neighboring heading.

For $g\geq3$, the supplied summary attributes to Korkmaz the statement that every homomorphism $\operatorname{Mod}(S_g)\to\operatorname{Out}(F_n)$ is trivial when $n\leq2g-1$. At $n=2g$, it reports a reduction rather than a classification. Let
\[
 q:\operatorname{Out}(F_{2g})\to\operatorname{GL}_{2g}(\Z)
\]
be the action on free-group abelianization, let $\sigma_g:\operatorname{Mod}(S_g)\to\operatorname{Sp}_{2g}(\Z)$ be the standard action on surface homology, and let $\mathcal I_g=\ker\sigma_g$ be Torelli. For every nontrivial $\rho:\operatorname{Mod}(S_g)\to\operatorname{Out}(F_{2g})$, the aid states that the complex representation $q\circ\rho$ is conjugate in $\operatorname{GL}_{2g}(\C)$ to $\sigma_g$. Its stated consequences are
\[
 \ker\rho\subseteq\mathcal I_g,\qquad\ker\rho\text{ torsion-free},
 \qquad\rho|_H\text{ injective for every finite }H\leq\operatorname{Mod}(S_g),
 \qquad|\rho(\operatorname{Mod}(S_g))|=\infty.
\]
This is a complex-conjugacy homological constraint, not a claimed integral conjugacy or proof that every such $\rho$ is injective.

The aid separately reports finite-image maps in larger target rank. Reduction modulo two gives an action on the nonzero vectors of $\F_2^{2g}$, a set of size $D=2^{2g}-1$. Permuting the corresponding free basis gives a map into $\operatorname{Out}(F_D)$ with image isomorphic to $\operatorname{Sp}_{2g}(\F_2)$. These do not contradict the infinite-image conclusion at rank $2g$. It also mentions embeddings in genera one and two without supplying their target ranks or formulas; those are not guessed. The forward classification, the first unresolved rank, and the reverse-direction research questions remain separate.
\subsection{Short English statement}
The source asks about maps from closed-surface mapping class groups to free-group outer automorphism groups, with reverse-direction questions in the same extract. The supplied partial result kills low target ranks and constrains any nontrivial map at rank $2g$ to have the standard symplectic homological shadow. It does not turn that shadow into a complete classification or settle the reverse questions.
\subsection{Sources}
\begin{itemize}
\item Local UnsolvedMath registry, permanent ID 20001650, exact source code catalogue entry; canonical \texttt{aim-geometric-group-theory-notes.json}, \texttt{notes[158]}, source item 4.1, tag ``problem'', workshop \emph{The geometry of the outer automorphism group of a free group}; \url{https://aimath.org/WWN/outerauto/outerauto.pdf}.
\item Original reverse-direction attribution: Brendle and Hamidi-Tehrani for exclusion of poison subgroups from mapping class groups. The summary names Korkmaz for the low-rank forward theorem, but supplies no titles, years, theorem numbers or complete bibliographic entries for either input. No proposer is identified; original remarks and literature field are empty.
\item Supplied metadata: PARTIAL-PROGRESS, registry status ``partially\_solved'', statement status ``exact''; attempt 1, result ``partial\_result'', problem status at run ``partially\_solved'', recorded rigor ``proved'', novelty confidence ``low'', candidate type ``reduction''. The machine-generated aid is not peer reviewed and explicitly requires independent verification before citation. No external source was visited.
\end{itemize}
\end{document}
